\subsection{MAGNUS GROUP}

For all \(a = (a_n)_{n \geq 0}\) in \(\hat{\mathbf{A}}(\mathbf{X})\) , the element \(a_0\) of \(\mathbf{K}\) will be called the constant term of \(a\) and denoted by \(\varepsilon(a)\) . Formula (1) shows that \(\varepsilon\) is an algebra homomorphism of \(\hat{\mathbf{A}}(\mathbf{X})\) into \(\mathbf{K}\) .

Lemma 1. For an element \(a\) of \(\hat{\mathbf{A}}(\mathbf{X})\) to be invertible, it is necessary and sufficient that its constant term be invertible in \(\mathbf{K}\) .

If \(a\) is invertible in \(\hat{\mathbf{A}}(\mathbf{X})\) , \(\varepsilon(a)\) is invertible in \(\mathbf{K}\) . Conversely, if \(\varepsilon(a)\) is invertible in \(\mathbf{K}\) , there exists \(u \in \hat{\mathbf{A}}_1(\mathbf{X})\) such that \(a = \varepsilon(a)(1 - u)\) ; we write \(b = \left(\sum_{n \geqslant 0} u^n\right)\varepsilon(a)^{-1}\) . Then \(ab = ba = 1\) and \(a\) is invertible.

The set of elements of \(\hat{\mathbf{A}}(\mathbf{X})\) of constant term 1 is therefore a subgroup of the multiplicative monoid \(\hat{\mathbf{A}}(\mathbf{X})\) , called the Magnus group constructed on \(\mathbf{X}\) (relative to \(\mathbf{K}\) ). In this chapter it will be denoted by \(\Gamma(\mathbf{X})\) or simply \(\Gamma\) . For every integer \(n \geqslant 1\) , we denote by \(\Gamma_n\) the set of \(a \in \Gamma\) such that \(\omega(a - 1) \geqslant n\) . By Proposition 2 of §4, no. 5, the sequence \((\Gamma_n)_{n \geqslant 1}\) is an integral central filtration on \(\Gamma\) .

